\chapter{Struktur Lewis, Resonansi, dan VSEPR}\label{ch:b1:lewis-resonance-vsepr}

Molekul ozon, \ce{O3}, terdiri atas tiga atom oksigen dalam rantai yang
bengkok. Gambarlah molekul itu dengan aturan-aturan Buku 1, maka satu
ikatan menjadi rangkap dua dan yang lain tunggal: satu ikatan semestinya
pendek, yang lain panjang. Namun, spektroskopi gelombang mikro mendapati
kedua ikatan \ce{O-O} ozon tepat sama panjang, yaitu \qty{127.8}{pm}, di
antara ikatan rangkap dua \ce{O2} (\qty{120.8}{pm}) dan ikatan tunggal
hidrogen peroksida (\qty{147.5}{pm}). Tidak ada satu gambar pun yang
benar; dua gambar bersama-samalah yang benar. Bab ini mempertajam model
Lewis --- \omterm{def:b1:lewis-resonance-vsepr:formal-charge}{muatan formal}, oktet yang tidak lengkap atau terlampaui,
resonansi antara beberapa struktur --- lalu memakainya untuk meramalkan
bentuk molekul dan \omterm{def:b1:lewis-resonance-vsepr:dipole}{momen dipol} yang mengikutinya.
% ledger: geo:O3.rOO, re:O2, geo:H2O2.rOO

\begin{recall}
Buku 1 (kelas 10) memerikan \omterm{def:b1:lewis-resonance-vsepr:lewis}{ikatan kovalen} sebagai pasangan elektron yang
dipakai bersama, menggambar \omterm{def:b1:lewis-resonance-vsepr:lewis}{struktur Lewis} dengan \omterm{def:b1:lewis-resonance-vsepr:lewis}{pasangan elektron
ikatan} dan \omterm{def:b1:lewis-resonance-vsepr:lewis}{pasangan elektron bebas}, dan meramalkan empat bentuk molekul
(linear, bengkok, segitiga datar, tetrahedral). \omterm{def:b1:periodicity:electronegativity}{Keelektronegatifan} dan
selisihnya dikuantitatifkan pada \cref{ch:b1:periodicity}; \omterm{def:b1:quantum-numbers:configuration}{elektron
valensi} sebuah atom dibaca dari konfigurasinya
(\cref{def:b1:quantum-numbers:configuration}).
\end{recall}

\section{Struktur Lewis dan muatan formal}

\begin{definition}[Ikatan kovalen, struktur Lewis]\label{def:b1:lewis-resonance-vsepr:lewis}
\emph{Ikatan kovalen}\index{ikatan kovalen} adalah sepasang elektron yang
dipakai bersama oleh dua atom, yaitu \emph{pasangan elektron
ikatan}\index{pasangan elektron ikatan}; dua atau tiga pasang membentuk
ikatan rangkap dua atau rangkap tiga. Sepasang \omterm{def:b1:quantum-numbers:configuration}{elektron valensi} yang
dimiliki satu atom saja adalah \emph{pasangan elektron
bebas}\index{pasangan elektron bebas}. Sebuah \emph{struktur
Lewis}\index{struktur Lewis} memperlihatkan setiap \omterm{def:b1:quantum-numbers:configuration}{elektron valensi}
sebuah molekul atau ion, sebagai ikatan (garis antaratom) dan pasangan
elektron bebas (garis atau sepasang titik pada sebuah atom). Menurut
\emph{aturan oktet}\index{aturan oktet}, atom-atom \omterm{def:b1:periodicity:table}{periode} 2 (C, N, O,
F) dalam struktur yang stabil dikelilingi empat pasang, yaitu delapan
elektron; menurut \emph{aturan duplet}\index{aturan duplet}, hidrogen
dikelilingi satu pasang.
\end{definition}

\begin{definition}[Muatan formal]\label{def:b1:lewis-resonance-vsepr:formal-charge}
\emph{Muatan formal}\index{muatan formal} sebuah atom dalam \omterm{def:b1:lewis-resonance-vsepr:lewis}{struktur
Lewis} adalah
\[
  q_F = N_v - N_{\text{bebas}} - \tfrac12 N_{\text{ikatan}},
\]
dengan $N_v$ banyaknya \omterm{def:b1:quantum-numbers:configuration}{elektron valensi} atom bebas, $N_{\text{bebas}}$
banyaknya elektron pada pasangan-pasangan elektron bebasnya, dan
$N_{\text{ikatan}}$ banyaknya elektron pada ikatan-ikatan yang
dibentuknya. Muatan ini ditulis $\oplus$ atau $\ominus$ di samping
atomnya jika tidak sama dengan nol.
\end{definition}

\begin{proposition}[Jumlah muatan formal sama dengan muatan spesies]\label{prop:b1:lewis-resonance-vsepr:formal-sum}
Jumlah \omterm{def:b1:lewis-resonance-vsepr:formal-charge}{muatan formal} atom-atom sebuah \omterm{def:b1:lewis-resonance-vsepr:lewis}{struktur Lewis} sama dengan muatan
total molekul atau ion itu.
\end{proposition}

\begin{proof}
Menjumlahkan $q_F$ untuk semua atom menghasilkan
$\sum N_v - (\text{elektron pada pasangan bebas} + \text{elektron pada ikatan})$,
karena elektron setiap ikatan dihitung separuh pada masing-masing kedua
atomnya. Isi kurung itu adalah banyaknya \omterm{def:b1:quantum-numbers:configuration}{elektron valensi} yang digambar
dalam struktur, yang sama dengan $\sum N_v$ dikurangi muatan $z$ spesies
itu. Karena itu $\sum q_F = z$.
\end{proof}

\begin{method}[Menggambar struktur Lewis]\label{met:b1:lewis-resonance-vsepr:lewis}
Untuk menggambar \omterm{def:b1:lewis-resonance-vsepr:lewis}{struktur Lewis} sebuah molekul atau ion:
\begin{enumerate}
  \item hitung \omterm{def:b1:quantum-numbers:configuration}{elektron valensi} $N = \sum N_v - z$ dan banyaknya pasangan
        $N/2$;
  \item gambarlah kerangkanya: atom yang paling tidak elektronegatif
        (bukan H) biasanya menjadi atom pusat; hidrogen dan fluorin selalu
        berada di ujung;
  \item lengkapi oktet atom-atom ujung dengan \omterm{def:b1:lewis-resonance-vsepr:lewis}{pasangan elektron bebas},
        lalu tempatkan pasangan yang tersisa pada atom pusat;
  \item jika atom pusat belum mencapai oktet, ubahlah \omterm{def:b1:lewis-resonance-vsepr:lewis}{pasangan elektron
        bebas} atom-atom tetangganya menjadi ikatan rangkap;
  \item hitunglah \omterm{def:b1:lewis-resonance-vsepr:formal-charge}{muatan formalnya}; di antara struktur-struktur yang
        mungkin, pilihlah yang \omterm{def:b1:lewis-resonance-vsepr:formal-charge}{muatan formalnya} paling sedikit, dengan
        muatan negatif pada atom-atom yang paling elektronegatif.
\end{enumerate}
\end{method}

\begin{example}[Asam nitrat]\label{ex:b1:lewis-resonance-vsepr:hno3}
\ce{HNO3}: $N = 1 + 5 + 3 \times 6 = 24$ elektron, 12 pasang. Kerangka:
nitrogen di pusat, terikat pada tiga oksigen, salah satunya membawa
hidrogen. Melengkapi oktet dengan ikatan tunggal menyisakan enam
elektron pada nitrogen; satu \omterm{def:b1:lewis-resonance-vsepr:lewis}{pasangan elektron bebas} oksigen menjadi
ikatan \ce{N=O}. \omterm{def:b1:lewis-resonance-vsepr:formal-charge}{Muatan formal}: nitrogen $5 - 0 - 4 = +1$; oksigen
berikatan tunggal tanpa hidrogen $6 - 6 - 1 = -1$; yang lain 0.
Jumlahnya 0, sebagaimana dituntut
\cref{prop:b1:lewis-resonance-vsepr:formal-sum}.
\end{example}

\begin{omfigure}
\setchemfig{atom sep=2.4em}
\begin{tikzpicture}[font=\small]
  \node at (0,0) {\large\chemfig{H-\charge{90=\|,270=\|}{O}-\charge{90:2pt=$\scriptstyle\oplus$}{N}(=[:45]\charge{0=\|,90=\|}{O})-[:-45]\charge{0=\|,240=\|,300=\|,45:3pt=$\scriptstyle\ominus$}{O}}};
  \node[below] at (0,-1.5) {asam nitrat, \ce{HNO3}};
  \node at (5.2,0) {\large\chemfig{\charge{180=\|}{N}~\charge{0=\|}{N}}};
  \node[below] at (5.2,-1.5) {dinitrogen, \ce{N2}};
  \node at (9.6,0) {\large\chemfig{\charge{135=\|,225=\|,90:2pt=$\scriptstyle\ominus$}{N}=\charge{90:2pt=$\scriptstyle\oplus$}{N}=\charge{45=\|,315=\|}{O}}};
  \node[below] at (9.6,-1.5) {dinitrogen monoksida, \ce{N2O}};
\end{tikzpicture}

{\small \omterm{def:b1:lewis-resonance-vsepr:lewis}{Struktur Lewis} dengan \omterm{def:b1:lewis-resonance-vsepr:lewis}{pasangan elektron bebas} dan \omterm{def:b1:lewis-resonance-vsepr:formal-charge}{muatan formal}.}
\end{omfigure}

\section{Melampaui oktet}

\begin{definition}[Molekul kekurangan elektron dan molekul hipervalen]\label{def:b1:lewis-resonance-vsepr:hypervalent}
Molekul yang salah satu atomnya mempunyai kurang dari delapan \omterm{def:b1:quantum-numbers:configuration}{elektron
valensi} disebut \emph{kekurangan elektron}\index{kekurangan elektron}:
dalam \ce{BF3}, boron mempunyai enam. Molekul yang salah satu atomnya,
dari \omterm{def:b1:periodicity:table}{periode} 3 atau sesudahnya, dikelilingi lebih dari empat pasang
disebut \emph{hipervalen}\index{hipervalen}: dalam \ce{PCl5} fosfor
membentuk lima ikatan, dalam \ce{SF6} belerang enam. Molekul dengan
jumlah elektron ganjil, seperti \ce{NO} dan \ce{NO2}, tidak dapat
memenuhi \omterm{def:b1:lewis-resonance-vsepr:lewis}{aturan oktet} pada setiap atom: satu elektron tidak
berpasangan.
\end{definition}

\begin{remark}[Mengapa periode 3 dan bukan periode 2]\label{rem:b1:lewis-resonance-vsepr:hypervalent}
Nitrogen membentuk \ce{NF3} tetapi tidak pernah \ce{NF5}, sedangkan fosfor
membentuk \ce{PF3} dan juga \ce{PF5}. Atom \omterm{def:b1:periodicity:table}{periode} 2 terlalu kecil untuk
berikatan dengan lebih dari empat tetangga; atom \omterm{def:b1:periodicity:table}{periode} 3 cukup besar
untuk menampung lima atau enam. Gambar \ce{SO4^{2-}} atau \ce{H2SO4}
dengan ikatan rangkap pada belerang (12 elektron di sekitar S) maupun
dengan ikatan tunggal dan \omterm{def:b1:lewis-resonance-vsepr:formal-charge}{muatan formal} (sebuah oktet) sama-sama
dijumpai; gambar kedua lebih mencerminkan sebaran muatannya, dan
keduanya menghasilkan geometri yang benar.
\end{remark}

\begin{definition}[Asam Lewis, basa Lewis, ikatan datif]\label{def:b1:lewis-resonance-vsepr:lewis-acid}
\emph{Asam Lewis}\index{asam Lewis} adalah spesies yang mampu menerima
sepasang elektron ke dalam tempat kosong pada kulit valensinya (\ce{BF3},
\ce{AlCl3}, \ce{H+}, kation logam); \emph{basa Lewis}\index{basa Lewis}
adalah spesies yang mempunyai \omterm{def:b1:lewis-resonance-vsepr:lewis}{pasangan elektron bebas} yang dapat
dipakainya bersama (\ce{NH3}, \ce{H2O}, \ce{F-}). Jika basa memberikan
pasangan itu kepada asam, ikatan yang terbentuk adalah \emph{ikatan
datif}\index{ikatan datif}: \omterm{def:b1:lewis-resonance-vsepr:lewis}{ikatan kovalen} yang kedua elektronnya berasal
dari satu pihak.
\end{definition}

\begin{example}[Amonia dan boron trifluorida]\label{ex:b1:lewis-resonance-vsepr:adduct}
\ce{NH3 + BF3 -> H3NBF3}. Dalam aduknya, boron melengkapi oktetnya;
nitrogen, yang kini memakai bersama \omterm{def:b1:lewis-resonance-vsepr:lewis}{pasangan elektron bebasnya}, membawa
\omterm{def:b1:lewis-resonance-vsepr:formal-charge}{muatan formal} $+1$ dan boron $-1$. Setelah terbentuk, ikatan \ce{N-B}
adalah \omterm{def:b1:lewis-resonance-vsepr:lewis}{ikatan kovalen} biasa.
\end{example}

\begin{omfigure}
\setchemfig{atom sep=2.1em}
\begin{tikzpicture}[font=\small]
  \node at (0,0) {\large\chemfig{H-\charge{90=\|}{N}(-[:240]H)(-[:300]H)}};
  \node at (1.6,0) {$+$};
  \node at (3.1,0) {\large\chemfig{B(-[:90]F)(-[:210]F)(-[:330]F)}};
  \draw[-{Stealth}, thick] (4.3,0) -- (5.3,0);
  \node at (7.4,0) {\large\chemfig{H-\charge{45:1pt=$\scriptstyle\oplus$}{N}(-[:90]H)(-[:270]H)-\charge{45:1pt=$\scriptstyle\ominus$}{B}(-[:90]F)(-[:270]F)-F}};
\end{tikzpicture}

{\small \omterm{def:b1:lewis-resonance-vsepr:lewis-acid}{Ikatan datif}: \omterm{def:b1:lewis-resonance-vsepr:lewis}{pasangan elektron bebas} amonia (\omterm{def:b1:lewis-resonance-vsepr:lewis-acid}{basa Lewis}) mengisi
tempat kosong boron trifluorida (\omterm{def:b1:lewis-resonance-vsepr:lewis-acid}{asam Lewis}). \omterm{def:b1:lewis-resonance-vsepr:lewis}{Pasangan elektron bebas}
fluorin tidak digambar.}
\end{omfigure}

\section{Resonansi}

\begin{definition}[Struktur resonansi, hibrida resonansi]\label{def:b1:lewis-resonance-vsepr:resonance}
Jika elektron-elektron sebuah molekul dapat ditempatkan dalam beberapa
\omterm{def:b1:lewis-resonance-vsepr:lewis}{struktur Lewis} yang hanya berbeda dalam letak ikatan $\pi$ dan \omterm{def:b1:lewis-resonance-vsepr:lewis}{pasangan
elektron bebasnya} --- sementara inti-inti tetap di tempatnya --- setiap
struktur itu adalah sebuah
\emph{struktur resonansi}\index{struktur resonansi} (atau bentuk mesomeri), yang dihubungkan dengan struktur
lainnya oleh panah berkepala dua $\leftrightarrow$. Molekul yang
sebenarnya bukan salah satu di antaranya: molekul itu adalah
\emph{hibrida resonansi}\index{hibrida resonansi}, satu struktur tunggal
yang sebaran elektronnya merupakan rata-rata berbobot sebaran
struktur-struktur itu.
\end{definition}

\begin{remark}[Sebuah panah, bukan kesetimbangan]\label{rem:b1:lewis-resonance-vsepr:arrow}
Panah $\leftrightarrow$ tidak memerikan sebuah reaksi: ozon tidak beralih
dari satu struktur ke struktur yang lain. Hibrida itu adalah satu
molekul, yang dipaparkan dengan beberapa gambar karena satu gambar tidak
cukup. Jangan pernah memakai $\rightleftharpoons$ di antara
struktur-struktur resonansi.
\end{remark}

\begin{method}[Menulis dan menimbang struktur resonansi]\label{met:b1:lewis-resonance-vsepr:resonance}
Untuk menemukan struktur-struktur resonansi sebuah spesies:
\begin{enumerate}
  \item mulailah dari satu \omterm{def:b1:lewis-resonance-vsepr:lewis}{struktur Lewis}; carilah \omterm{def:b1:lewis-resonance-vsepr:lewis}{pasangan elektron
        bebas} atau ikatan $\pi$ yang bersebelahan dengan ikatan $\pi$,
        tempat kosong, atau muatan positif;
  \item pindahkan pasangan-pasangan elektron dengan panah lengkung
        (\omterm{def:b1:lewis-resonance-vsepr:lewis}{pasangan bebas} menjadi ikatan $\pi$, ikatan $\pi$ menjadi
        \omterm{def:b1:lewis-resonance-vsepr:lewis}{pasangan bebas}), tanpa pernah memindahkan inti dan tanpa pernah
        melampaui oktet atom \omterm{def:b1:periodicity:table}{periode} 2;
  \item hitung ulang \omterm{def:b1:lewis-resonance-vsepr:formal-charge}{muatan formalnya};
  \item timbanglah struktur-struktur itu: yang paling penting mematuhi
        \omterm{def:b1:lewis-resonance-vsepr:lewis}{aturan oktet}, mempunyai \omterm{def:b1:lewis-resonance-vsepr:formal-charge}{muatan formal} paling sedikit, dan
        menempatkan muatan negatif pada atom-atom yang paling
        elektronegatif; struktur-struktur yang ekuivalen berbobot sama.
\end{enumerate}
\end{method}

\begin{example}[Ozon, nitrat, benzena]\label{ex:b1:lewis-resonance-vsepr:examples}
Ozon mempunyai dua struktur yang ekuivalen: setiap ikatan \ce{O-O}
rangkap dua pada struktur yang satu dan tunggal pada struktur yang lain,
sehingga keduanya mempunyai orde ikatan $1.5$ dan panjang yang sama,
\qty{127.8}{pm}, di antara \ce{O=O} dan \ce{O-O}. Ion nitrat \ce{NO3-}
mempunyai tiga struktur yang ekuivalen; setiap ikatan \ce{N-O} berorde
$4/3$. Benzena, \ce{C6H6}, mempunyai dua struktur Kekulé; keenam ikatan
\ce{C-C}-nya sama panjang, \qty{139.7}{pm}, di antara ikatan tunggal
etana (\qty{153.6}{pm}) dan ikatan rangkap dua etena (\qty{133.9}{pm}).
% ledger: geo:O3.rOO, geo:C6H6.rCC, geo:C2H6.rCC, geo:C2H4.rCC
\end{example}

\begin{omfigure}
\vspace*{10pt}
\large
\setchemfig{atom sep=2.4em}
\schemestart
\chemfig{\charge{135=\|,225=\|}{O}=[:-30]\charge{270=\|,90:2pt=$\scriptstyle\oplus$}{O}-[:30]\charge{90=\|,0=\|,270=\|,45:5pt=$\scriptstyle\ominus$}{O}}
\arrow{<->}
\chemfig{\charge{90=\|,180=\|,270=\|,135:5pt=$\scriptstyle\ominus$}{O}-[:-30]\charge{270=\|,90:2pt=$\scriptstyle\oplus$}{O}=[:30]\charge{45=\|,315=\|}{O}}
\schemestop

\medskip
\schemestart
\chemfig{*6(=-=-=-)}
\arrow{<->}
\chemfig{*6(-=-=-=)}
\schemestop
\hspace{3em}
\chemfig{\charge{270:2pt=$\scriptstyle\oplus$}{N}(=[:90]\charge{45=\|,135=\|}{O})(-[:210]\charge{150=\|,240=\|,300=\|,90:2pt=$\scriptstyle\ominus$}{O})-[:330]\charge{30=\|,300=\|,240=\|,90:2pt=$\scriptstyle\ominus$}{O}}
\normalsize
\par\vspace{14pt}
{\small Resonansi pada ozon (satu \omterm{def:b1:lewis-resonance-vsepr:lewis}{pasangan elektron bebas} oksigen kanan
menjadi ikatan $\pi$ sementara ikatan $\pi$ di kiri menjadi \omterm{def:b1:lewis-resonance-vsepr:lewis}{pasangan
bebas}) dan pada benzena, serta salah satu dari tiga struktur ekuivalen
ion nitrat, yang muatannya terbagi pada ketiga oksigen.}
\end{omfigure}

\begin{definition}[Ikatan $\sigma$, ikatan $\pi$]\label{def:b1:lewis-resonance-vsepr:sigma-pi}
Ikatan tunggal adalah \emph{ikatan $\sigma$}\index{ikatan sigma@ikatan $\sigma$}:
pasangan elektronnya terletak di sepanjang sumbu yang menghubungkan kedua
inti, hasil tumpang tindih ujung-ke-ujung dua orbital. Ikatan kedua dan
ketiga pada ikatan rangkap adalah \emph{ikatan $\pi$}\index{ikatan pi@ikatan $\pi$}:
pasangan-pasangannya terletak di kedua sisi sumbu, hasil tumpang tindih
sisi-ke-sisi dua orbital p yang tegak lurus pada sumbu itu. Ikatan
rangkap dua terdiri atas satu ikatan $\sigma$ dan satu ikatan $\pi$;
ikatan rangkap tiga atas satu $\sigma$ dan dua $\pi$.
\end{definition}

\begin{omfigure}
\begin{tikzpicture}[font=\small]
  % sigma: two lobes overlapping head-on
  \fill[omDef!35, opacity=0.8] (0,0) ellipse (0.9 and 0.42);
  \fill[omProp!35, opacity=0.8] (1.1,0) ellipse (0.9 and 0.42);
  \fill (0.1,0) circle (1.6pt); \fill (1.0,0) circle (1.6pt);
  \draw[dashed, black!50] (-1.2,0) -- (2.3,0);
  \node[below, align=center] at (0.55,-1.05) {$\sigma$: ujung-ke-ujung,\\pada sumbu};
  % pi: two p orbitals side by side
  \begin{scope}[xshift=5.2cm]
    \foreach \x/\c in {0/omDef, 1.2/omProp} {
      \fill[\c!35, opacity=0.85] (\x,0.45) ellipse (0.32 and 0.45);
      \fill[\c!15, opacity=0.85] (\x,-0.45) ellipse (0.32 and 0.45);
      \draw[\c!70!black] (\x,0.45) ellipse (0.32 and 0.45) (\x,-0.45) ellipse (0.32 and 0.45);
      \fill (\x,0) circle (1.6pt);
    }
    \draw[dashed, black!50] (-0.7,0) -- (1.9,0);
    \draw[omThm, thick, densely dotted] (0.25,0.75) -- (0.95,0.75) (0.25,-0.75) -- (0.95,-0.75);
    \node[below, align=center] at (0.6,-1.05) {$\pi$: sisi-ke-sisi,\\di atas dan di bawah};
  \end{scope}
\end{tikzpicture}

{\small Dua jenis tumpang tindih. Ikatan $\sigma$ memungkinkan rotasi
mengelilingi sumbunya; ikatan $\pi$ tidak, itulah sebabnya ikatan
rangkap dua kaku (\cref{ch:b1:stereochemistry-in-depth}).}
\end{omfigure}

\section{Model VSEPR}

\begin{definition}[Model VSEPR, bilangan sterik]\label{def:b1:lewis-resonance-vsepr:vsepr}
Dalam \emph{model VSEPR}\index{model VSEPR} (tolakan pasangan elektron
kulit valensi), pasangan-pasangan elektron di sekitar atom pusat A saling
tolak dan mengambil susunan yang membuat mereka sejauh mungkin satu sama
lain. Sebuah molekul ditulis \ce{AX_nE_m}, dengan $n$ banyaknya atom
yang terikat pada A (ikatan rangkap dihitung sekali) dan $m$ banyaknya
\omterm{def:b1:lewis-resonance-vsepr:lewis}{pasangan elektron bebas} A. \emph{Bilangan sterik}\index{bilangan sterik}
$n + m$ menentukan susunan pasangan-pasangan itu: 2 linear, 3 segitiga
datar, 4 tetrahedral, 5 bipiramida segitiga, 6 oktahedral. Bentuk
molekul adalah susunan atom-atomnya saja.
\end{definition}

\begin{proposition}[Geometri VSEPR]\label{prop:b1:lewis-resonance-vsepr:geometries}
\omterm{def:b1:lewis-resonance-vsepr:vsepr}{Model VSEPR} meramalkan bentuk dan sudut ideal pada tabel di bawah ini;
sudut-sudut terukur mendekatinya.
\begin{center}
\small
\begin{tabular}{llllc}
\hline
tipe & susunan pasangan & bentuk & contoh & sudut terukur \\
\hline
\ce{AX2} & linear & linear & \ce{CO2} & $180^\circ$ \\
\ce{AX3} & segitiga datar & segitiga datar & \ce{BF3} & $120^\circ$ \\
\ce{AX2E} & segitiga datar & bengkok & \ce{SO2} & $119.5^\circ$ \\
\ce{AX4} & tetrahedral & tetrahedral & \ce{CH4} & $109.5^\circ$ \\
\ce{AX3E} & tetrahedral & piramida segitiga & \ce{NH3} & $106.7^\circ$ \\
\ce{AX2E2} & tetrahedral & bengkok & \ce{H2O} & $104.5^\circ$ \\
\ce{AX5} & bipiramida segitiga & bipiramida segitiga & \ce{PCl5} & $90^\circ$, $120^\circ$ \\
\ce{AX4E} & bipiramida segitiga & jungkat-jungkit & \ce{SF4} & $101.6^\circ$, $173.1^\circ$ \\
\ce{AX3E2} & bipiramida segitiga & bentuk T & \ce{ClF3} & $87.5^\circ$ \\
\ce{AX2E3} & bipiramida segitiga & linear & \ce{XeF2} & $180^\circ$ \\
\ce{AX6} & oktahedral & oktahedral & \ce{SF6} & $90^\circ$ \\
\ce{AX5E} & oktahedral & piramida segi empat & \ce{BrF5} & --- \\
\ce{AX4E2} & oktahedral & segi empat datar & \ce{XeF4} & --- \\
\hline
\end{tabular}
\end{center}
% ledger: geo:CO2.aOCO, geo:BF3.aFBF, geo:SO2.aOSO, geo:CH4.aHCH,
% geo:NH3.aHNH, geo:H2O.aHOH, geo:SF6.aFSF, geo:SF4.aFSF2, geo:SF4.aFSF,
% geo:ClF3.aFClF, geo:XeF2.aFXeF
\end{proposition}
\admitted

\begin{proposition}[Pasangan elektron bebas menyempitkan sudut]\label{prop:b1:lewis-resonance-vsepr:lone-pair-angles}
\omterm{def:b1:lewis-resonance-vsepr:lewis}{Pasangan elektron bebas}, yang hanya diikat oleh satu inti, lebih
menyebar daripada \omterm{def:b1:lewis-resonance-vsepr:lewis}{pasangan elektron ikatan} dan menolak lebih kuat:
tolakan berkurang menurut urutan \omterm{def:b1:lewis-resonance-vsepr:lewis}{pasangan bebas}/\omterm{def:b1:lewis-resonance-vsepr:lewis}{pasangan bebas} $>$
\omterm{def:b1:lewis-resonance-vsepr:lewis}{pasangan bebas}/\omterm{def:b1:lewis-resonance-vsepr:lewis}{pasangan ikatan} $>$ \omterm{def:b1:lewis-resonance-vsepr:lewis}{pasangan ikatan}/\omterm{def:b1:lewis-resonance-vsepr:lewis}{pasangan ikatan}.
Karena itu sudut antarikatan menyempit ketika \omterm{def:b1:lewis-resonance-vsepr:lewis}{pasangan elektron bebas}
bertambah ($109.5^\circ$ pada \ce{CH4}, $106.7^\circ$ pada \ce{NH3},
$104.5^\circ$ pada \ce{H2O}), dan dalam bipiramida segitiga \omterm{def:b1:lewis-resonance-vsepr:lewis}{pasangan
elektron bebas} menempati posisi ekuatorial, tempat pasangan itu hanya
mempunyai dua tetangga pada $90^\circ$.
% ledger: geo:CH4.aHCH, geo:NH3.aHNH, geo:H2O.aHOH
\end{proposition}
\admitted

\begin{omfigure}
\setchemfig{atom sep=2.6em}
\renewcommand{\arraystretch}{1.2}
\begin{tabular}{ccccc}
\chemfig{O=C=O} &
\chemfig{B(-[:90]F)(-[:210]F)(-[:330]F)} &
\chemfig{C(-[:90]H)(-[:200]H)(<[:280]H)(<:[:335]H)} &
\chemfig{P(-[:90]Cl)(-[:270]Cl)(-[:180]Cl)(<[:320]Cl)(<:[:30]Cl)} &
\chemfig{S(-[:90]F)(-[:270]F)(-[:180]F)(-[:0]F)(<[:225]F)(<:[:45]F)} \\
\small \ce{AX2}, \ce{CO2} & \small \ce{AX3}, \ce{BF3} & \small \ce{AX4}, \ce{CH4} &
\small \ce{AX5}, \ce{PCl5} & \small \ce{AX6}, \ce{SF6} \\[16pt]
\chemfig{\charge{90=\:}{S}(=[:210]O)(=[:330]O)} &
\chemfig{\charge{90=\:}{N}(-[:200]H)(<[:280]H)(<:[:335]H)} &
\chemfig{\charge{60=\:,120=\:}{O}(-[:215]H)(-[:325]H)} &
\chemfig{\charge{0=\:}{S}(-[:90]F)(-[:270]F)(<[:210]F)(<:[:150]F)} &
\chemfig{Xe(-[:0]F)(-[:90]F)(-[:180]F)(-[:270]F)} \\
\small \ce{AX2E}, \ce{SO2} & \small \ce{AX3E}, \ce{NH3} & \small \ce{AX2E2}, \ce{H2O} &
\small \ce{AX4E}, \ce{SF4} & \small \ce{AX4E2}, \ce{XeF4} \\
\end{tabular}

\medskip
{\small Bentuk-bentuk yang diramalkan VSEPR. Baji tebal mengarah ke
pembaca, baji bergaris menjauhi pembaca; ikatan biasa terletak pada
bidang halaman. \omterm{def:b1:lewis-resonance-vsepr:lewis}{Pasangan elektron bebas} atom pusat digambar sebagai
sepasang titik; \ce{XeF4} digambar tampak muka, dengan kedua \omterm{def:b1:lewis-resonance-vsepr:lewis}{pasangan
bebasnya} mengarah ke dan menjauhi pembaca.}
\end{omfigure}

\begin{method}[Meramalkan bentuk molekul]\label{met:b1:lewis-resonance-vsepr:vsepr}
Untuk meramalkan bentuk di sekitar sebuah atom pusat:
\begin{enumerate}
  \item gambarlah \omterm{def:b1:lewis-resonance-vsepr:lewis}{struktur Lewisnya} (\cref{met:b1:lewis-resonance-vsepr:lewis});
  \item hitunglah atom-atom yang terikat pada atom pusat ($n$) dan
        \omterm{def:b1:lewis-resonance-vsepr:lewis}{pasangan elektron bebasnya} ($m$); tuliskan \ce{AX_nE_m};
  \item bacalah susunannya dari $n + m$ dan bentuknya dari posisi
        $n$ atom itu;
  \item koreksilah sudut-sudut idealnya: \omterm{def:b1:lewis-resonance-vsepr:lewis}{pasangan elektron bebas} dan
        ikatan rangkap memerlukan ruang lebih banyak daripada ikatan
        tunggal.
\end{enumerate}
\end{method}

\begin{definition}[Hibridisasi]\label{def:b1:lewis-resonance-vsepr:hybridisation}
\emph{Hibridisasi}\index{hibridisasi} sebuah atom adalah deskripsi
ikatan-ikatannya dengan orbital-orbital yang mengarah sepanjang arah
yang diramalkan VSEPR: atom dengan \omterm{def:b1:lewis-resonance-vsepr:vsepr}{bilangan sterik} 4 disebut sp$^3$
(empat orbital ekuivalen pada $109.5^\circ$), dengan \omterm{def:b1:lewis-resonance-vsepr:vsepr}{bilangan sterik} 3
sp$^2$ (tiga orbital pada $120^\circ$ dalam satu bidang, dengan satu
orbital p tersisa yang tegak lurus pada bidang itu), dengan \omterm{def:b1:lewis-resonance-vsepr:vsepr}{bilangan
sterik} 2 sp (dua orbital pada $180^\circ$, dengan dua orbital p
tersisa). Orbital-orbital p yang tersisa membentuk ikatan-ikatan $\pi$.
\end{definition}

\begin{example}[Atom-atom karbon dalam kimia organik]\label{ex:b1:lewis-resonance-vsepr:carbon}
Dalam etana setiap karbon sp$^3$, tetrahedral; dalam etena sp$^2$,
keenam atomnya dalam satu bidang dengan sudut mendekati $120^\circ$
(\ce{H-C-H} terukur $117.6^\circ$); dalam etuna sp, keempat atomnya
pada satu garis. Panjang ikatannya menyusut dari 153.6 ke 133.9 ke
\qty{120.3}{pm} seiring ditambahkannya ikatan $\pi$.
% ledger: geo:C2H4.aHCH, geo:C2H6.rCC, geo:C2H4.rCC, geo:C2H2.rCC
\end{example}

\begin{remark}[Sebuah deskripsi, bukan sebab]\label{rem:b1:lewis-resonance-vsepr:hybrid-limit}
\omterm{def:b1:lewis-resonance-vsepr:hybridisation}{Hibridisasi} memerikan geometri yang sudah diketahui; \omterm{def:b1:lewis-resonance-vsepr:hybridisation}{hibridisasi} tidak
menjelaskannya, dan gagal untuk energi elektron, yang dijelaskan oleh
teori orbital molekul dalam jilid Tahun 2. Namun, \omterm{def:b1:lewis-resonance-vsepr:hybridisation}{hibridisasi} tetap
menjadi bahasa sehari-hari kimia organik: ``karbon sp$^2$'' berarti
karbon datar dengan satu orbital p yang bebas untuk ikatan $\pi$.
\end{remark}

\section{Molekul polar: momen dipol}

\begin{definition}[Dipol ikatan, momen dipol]\label{def:b1:lewis-resonance-vsepr:dipole}
Dalam ikatan antara atom-atom yang \omterm{def:b1:periodicity:electronegativity}{keelektronegatifannya} berbeda,
elektron-elektron condong ke atom yang lebih elektronegatif, yang membawa
muatan parsial $-\delta e$, sementara pasangannya membawa $+\delta e$
($0 < \delta < 1$). \emph{Dipol ikatan}\index{dipol ikatan} adalah vektor
$\vec\mu = \delta e\,
\vec d$, dengan norma $\delta e d$, yang mengarah
dari muatan parsial negatif ke muatan parsial positif, dengan $\vec d$
vektor di antara keduanya. \emph{Momen dipol}\index{momen dipol} sebuah
molekul adalah jumlah vektor dipol-dipol ikatannya (dan sumbangan
pasangan-pasangan elektron bebasnya); momen dipol dinyatakan dalam
coulomb meter atau dalam debye, $\qty{1}{\debye} =
\qty{3.336e-30}{C.m}$. Molekul dengan momen dipol yang tidak nol adalah
\emph{molekul polar}\index{molekul polar}.
% ledger: const:c (1 D = 1e-21/c C.m)
\end{definition}

\begin{remark}[Konvensi arah]\label{rem:b1:lewis-resonance-vsepr:convention}
Fisika menggambar vektor dipol dari muatan negatif ke muatan positif,
konvensi yang dipakai di sini. Buku-buku kimia yang lebih tua menggambar
panah berekor silang yang mengarah ke ujung negatif; hanya arahnya yang
berbeda.
\end{remark}

\begin{proposition}[Dipol molekul simetris]\label{prop:b1:lewis-resonance-vsepr:dipole-sum}
Molekul \ce{AX_n} tanpa \omterm{def:b1:lewis-resonance-vsepr:lewis}{pasangan elektron bebas} pada A (\ce{AX2} linear,
\ce{AX3} segitiga, \ce{AX4} tetrahedral, \dots) mempunyai \omterm{def:b1:lewis-resonance-vsepr:dipole}{momen dipol}
nol, betapa pun polarnya ikatan-ikatannya. Molekul \ce{AX2} bengkok
bersudut $\theta$ yang ikatan-ikatannya mempunyai dipol $\mu_b$ mempunyai
\omterm{def:b1:lewis-resonance-vsepr:dipole}{momen dipol}
\[
  \mu = 2\mu_b \cos\frac{\theta}{2} .
\]
\end{proposition}

\begin{proof}
Dalam bentuk-bentuk simetris, vektor-vektor ikatan sama normanya dan
arahnya tersusun simetris di sekitar A, sehingga jumlahnya nol (untuk
\ce{AX2} linear vektor-vektornya berlawanan; untuk \ce{AX3} pada
$120^\circ$ tiga vektor yang sama normanya berjumlah nol; untuk \ce{AX4}
jumlahnya invarian terhadap rotasi-rotasi tetrahedron, sehingga sama
dengan nol). Untuk molekul yang bengkok, setiap vektor ikatan membentuk
sudut $\theta/2$ dengan garis bagi sudutnya; komponen-komponen yang
tegak lurus pada garis bagi saling meniadakan dan komponen-komponen
sepanjang garis bagi saling menjumlah: $2\mu_b\cos(\theta/2)$.
\end{proof}

\begin{omfigure}
\begin{tikzpicture}[font=\small, >={Stealth}]
  % water
  \node[aO] (o) at (0,0) {};
  \node[aH] (h1) at (-0.95,-0.73) {};
  \node[aH] (h2) at (0.95,-0.73) {};
  \begin{scope}[on background layer]
    \draw[bond] (o.center) -- (h1.center); \draw[bond] (o.center) -- (h2.center);
  \end{scope}
  \draw[->, omProp, thick] (-0.25,0.15) -- (-0.95,-0.38);
  \draw[->, omProp, thick] (0.25,0.15) -- (0.95,-0.38);
  \draw[->, omThm, very thick] (0,-0.25) -- (0,-1.45);
  \node[omThm, right] at (0.05,-1.25) {$\vec\mu$};
  \node at (0,0.55) {$\delta^-$};
  \node at (-1.25,-1.05) {$\delta^+$}; \node at (1.25,-1.05) {$\delta^+$};
  \node[below] at (0,-1.7) {\ce{H2O}: $\mu = \qty{1.86}{\debye}$};
  % carbon dioxide
  \begin{scope}[xshift=5cm]
    \node[aO] (a) at (-1.1,0) {}; \node[aC] (c) at (0,0) {}; \node[aO] (b) at (1.1,0) {};
    \begin{scope}[on background layer]
      \foreach \d in {-1.6,1.6} {
        \draw[bond, line width=1.1pt] ([yshift=\d pt]a.center) -- ([yshift=\d pt]c.center);
        \draw[bond, line width=1.1pt] ([yshift=\d pt]c.center) -- ([yshift=\d pt]b.center);}
    \end{scope}
    \draw[->, omProp, thick] (-0.75,0.45) -- (-0.2,0.45);
    \draw[->, omProp, thick] (0.75,0.45) -- (0.2,0.45);
    \node[below, align=center] at (0,-1.7) {\ce{CO2}: dipol ikatan\\saling meniadakan, $\mu = 0$};
  \end{scope}
\end{tikzpicture}

{\small Dipol-dipol ikatan (jingga, dari muatan parsial negatif ke muatan
parsial positif) dan jumlahnya (merah). Pada air yang bengkok
dipol-dipol itu saling menjumlah; pada karbon dioksida yang linear
saling meniadakan.}
% ledger: dip:H2O
\end{omfigure}

\begin{example}[Molekul polar dan nonpolar]\label{ex:b1:lewis-resonance-vsepr:dipoles}
\ce{CO2}, \ce{BF3}, \ce{CH4}, dan \ce{CCl4} mempunyai \omterm{def:b1:lewis-resonance-vsepr:dipole}{momen dipol} nol.
\ce{H2O} (\qty{1.86}{\debye}), \ce{NH3} (\qty{1.48}{\debye}), dan
\ce{SO2} (\qty{1.63}{\debye}) bersifat polar. Sepanjang \ce{CH3Cl},
\ce{CH2Cl2}, \ce{CHCl3}, momennya turun, 1.87, 1.62, \qty{1.04}{\debye},
hingga lenyap pada \ce{CCl4}: dipol-dipol ikatan \ce{C-Cl} makin saling
meniadakan. Di antara hidrogen halida, momen itu mengikuti selisih
\omterm{def:b1:periodicity:electronegativity}{keelektronegatifan}: HF 1.83, HCl 1.09, HBr 0.83, HI \qty{0.45}{\debye}.
% ledger: dip:H2O, dip:NH3, dip:SO2, dip:CH3Cl, dip:CH2Cl2, dip:CHCl3,
% dip:CCl4, dip:HF, dip:HCl, dip:HBr, dip:HI
\end{example}

\section{Latihan}

\begin{exercise}[$\star$]\label{exo:b1:lewis-resonance-vsepr:1}
Gambarlah \omterm{def:b1:lewis-resonance-vsepr:lewis}{struktur Lewis} \ce{H2O}, \ce{NH3}, \ce{CO2}, \ce{HCN},
\ce{CH2O} (metanal), dan \ce{NH4+}, lalu tentukan \omterm{def:b1:lewis-resonance-vsepr:formal-charge}{muatan formalnya}.
\end{exercise}

\begin{exercise}[$\star$]\label{exo:b1:lewis-resonance-vsepr:2}
Tentukan tipe \ce{AX_nE_m}, bentuk, dan sudut perkiraan \ce{H2S},
\ce{PCl3}, \ce{BF3}, \ce{SiH4}, dan \ce{CO3^{2-}} (atom pusat ditulis
pertama dalam setiap rumus).
\end{exercise}

\begin{exercise}[$\star$]\label{exo:b1:lewis-resonance-vsepr:3}
Manakah di antara molekul-molekul berikut yang polar: \ce{CCl4},
\ce{CH2Cl2}, \ce{BF3}, \ce{NH3}, \ce{CO2}, \ce{SO2}? Berikan alasannya
dari bentuknya.
\end{exercise}

\begin{exercise}[$\star$]\label{exo:b1:lewis-resonance-vsepr:4}
Tentukan \omterm{def:b1:lewis-resonance-vsepr:lewis-acid}{asam Lewis} dan \omterm{def:b1:lewis-resonance-vsepr:lewis-acid}{basa Lewis} dalam \ce{H+ + H2O -> H3O+} dan dalam
\ce{Cu^{2+} + 4NH3 -> [Cu(NH3)4]^{2+}}. Ikatan manakah yang datif?
\end{exercise}

\begin{exercise}[$\star\star$]\label{exo:b1:lewis-resonance-vsepr:5}
Tuliskan dua \omterm{def:b1:lewis-resonance-vsepr:resonance}{struktur resonansi} ion etanoat \ce{CH3COO-}. Apa yang
diramalkan struktur-struktur itu tentang panjang kedua ikatan \ce{C-O}?
Bandingkan dengan asam metanoat \ce{HCOOH}, yang kedua ikatan
\ce{C-O}-nya berukuran 120.2 dan \qty{134.3}{pm}.
% ledger: geo:H2CO2.rCO, geo:H2CO2.rCO2
\end{exercise}

\begin{exercise}[$\star\star$]\label{exo:b1:lewis-resonance-vsepr:6}
Gambarlah \omterm{def:b1:lewis-resonance-vsepr:lewis}{struktur Lewis} ion sulfat \ce{SO4^{2-}} dengan empat ikatan
tunggal \ce{S-O}, hitunglah \omterm{def:b1:lewis-resonance-vsepr:formal-charge}{muatan formalnya}, dan ramalkan bentuknya.
Berapa banyak \omterm{def:b1:lewis-resonance-vsepr:resonance}{struktur resonansi} dengan dua ikatan \ce{S=O} dan tanpa
muatan pada belerang yang dapat ditulis?
\end{exercise}

\begin{exercise}[$\star\star$]\label{exo:b1:lewis-resonance-vsepr:7}
Gunakan VSEPR untuk meramalkan bentuk \ce{SF4}, \ce{ClF3}, \ce{XeF2},
dan \ce{XeF4}. Jelaskan mengapa pasangan-pasangan elektron bebas sebuah
bipiramida segitiga berada pada bidang ekuatorial.
\end{exercise}

\begin{exercise}[$\star\star$]\label{exo:b1:lewis-resonance-vsepr:8}
Sudut \ce{H-X-H} adalah $104.5^\circ$ pada \ce{H2O} dan $92.1^\circ$
pada \ce{H2S}; $106.7^\circ$ pada \ce{NH3} dan $93.3^\circ$ pada
\ce{PH3}. Kemukakan penjelasannya berdasarkan ukuran dan
\omterm{def:b1:periodicity:electronegativity}{keelektronegatifan} atom pusat.
% ledger: geo:H2O.aHOH, geo:H2S.aHSH, geo:NH3.aHNH, geo:PH3.aHPH
\end{exercise}

\begin{exercise}[$\star\star$]\label{exo:b1:lewis-resonance-vsepr:9}
Tentukan \omterm{def:b1:lewis-resonance-vsepr:hybridisation}{hibridisasi} setiap karbon dan nitrogen dalam \ce{CH3-C#N}
(etananitril) dan dalam \ce{CH2=CH-CH3}, serta banyaknya ikatan
$\sigma$ dan $\pi$ dalam setiap molekul.
\end{exercise}

\begin{exercise}[$\star\star\star$]\label{exo:b1:lewis-resonance-vsepr:10}
\omterm{def:b1:lewis-resonance-vsepr:dipole}{Momen dipol} \ce{NF3} hanya \qty{0.24}{\debye}, dibandingkan dengan
\qty{1.48}{\debye} untuk \ce{NH3}, padahal ikatan \ce{N-F} lebih polar
daripada ikatan \ce{N-H}. Dengan dipol-dipol ikatan dan \omterm{def:b1:lewis-resonance-vsepr:lewis}{pasangan
elektron bebas} nitrogen, jelaskan perbedaan itu.
% ledger: dip:NF3, dip:NH3
\end{exercise}

\begin{exercise}[$\star\star\star$]\label{exo:b1:lewis-resonance-vsepr:11}
Molekul \ce{SO2} mempunyai sudut $119.5^\circ$ dan \omterm{def:b1:lewis-resonance-vsepr:dipole}{momen dipol}
\qty{1.63}{\debye}. Hitunglah dipol satu ikatan \ce{S-O}; dengan panjang
ikatan \qty{143.2}{pm}, turunkan muatan parsial $\delta$ pada setiap
oksigen.
% ledger: geo:SO2.aOSO, geo:SO2.rSO, dip:SO2
\end{exercise}

\begin{exercise}[$\star\star\star$]\label{exo:b1:lewis-resonance-vsepr:12}
Gambarlah struktur-struktur Lewis terpenting dinitrogen monoksida
\ce{N2O} (linear, \ce{N-N-O}) dan hitunglah \omterm{def:b1:lewis-resonance-vsepr:formal-charge}{muatan formalnya}. Panjang
ikatan terukurnya \ce{N-N} \qty{112.8}{pm} dan \ce{N-O}
\qty{118.4}{pm}; dalam \ce{N2} ikatan rangkap tiganya berukuran
\qty{109.8}{pm}. Struktur manakah yang lebih berbobot?
% ledger: geo:N2O.rNN, geo:N2O.rNO, re:N2
\end{exercise}

\section{Soal: Molekul-Molekul Sambaran Petir}

\begin{problem}[{Soal akhir pekan --- spesies nitrogen dan oksigen yang
terbentuk ketika petir memanaskan udara: struktur Lewis, resonansi,
bentuk molekul, dan muatan parsial air}]
\label{pb:b1:lewis-resonance-vsepr:1}
Sambaran petir memanaskan udara hingga ribuan derajat: \ce{N2} dan
\ce{O2} menghasilkan \ce{NO}, lalu \ce{NO2}, \ce{N2O}, ozon, dan akhirnya
asam nitrat di dalam hujan. Data: $\qty{1}{\debye} = \qty{3.336e-30}{C.m}$,
$e =
\qty{1.602e-19}{C}$. Nilai-nilai terukur diberikan bila diperlukan.

\medskip
\noindent\textbf{Bagian I --- \omterm{def:b1:lewis-resonance-vsepr:lewis}{Struktur Lewis}.}
\begin{enumerate}
  \item Hitunglah \omterm{def:b1:quantum-numbers:configuration}{elektron valensi} \ce{N2}, \ce{NO}, \ce{NO2},
        \ce{N2O}, \ce{O3}, dan \ce{HNO3}.
  \item Gambarlah \omterm{def:b1:lewis-resonance-vsepr:lewis}{struktur Lewis} \ce{N2}.
  \item Gambarlah sebuah \omterm{def:b1:lewis-resonance-vsepr:lewis}{struktur Lewis} \ce{NO}. Mengapa \omterm{def:b1:lewis-resonance-vsepr:lewis}{aturan oktet}
        tidak dapat dipenuhi pada kedua atomnya?
  \item Gambarlah dua \omterm{def:b1:lewis-resonance-vsepr:lewis}{struktur Lewis} \ce{N2O}, dengan \ce{N=N=O} dan
        dengan \ce{N#N-O}, lalu tentukan \omterm{def:b1:lewis-resonance-vsepr:formal-charge}{muatan formalnya}.
  \item Gambarlah sebuah \omterm{def:b1:lewis-resonance-vsepr:lewis}{struktur Lewis} ozon dan tentukan \omterm{def:b1:lewis-resonance-vsepr:formal-charge}{muatan
        formalnya}.
  \item Gambarlah sebuah \omterm{def:b1:lewis-resonance-vsepr:lewis}{struktur Lewis} \ce{HNO3} dan tentukan \omterm{def:b1:lewis-resonance-vsepr:formal-charge}{muatan
        formalnya}.
  \item Manakah di antara keenam spesies itu yang mempunyai \omterm{def:b1:quantum-numbers:unpaired}{elektron tak
        berpasangan}?
\end{enumerate}

\medskip
\noindent\textbf{Bagian II --- Resonansi dan panjang ikatan.}
\begin{enumerate}[resume]
  \item Tuliskan dua \omterm{def:b1:lewis-resonance-vsepr:resonance}{struktur resonansi} ozon. Orde ikatan berapakah yang
        diramalkannya?
  \item Ikatan \ce{O-O} berukuran \qty{127.8}{pm} pada ozon,
        \qty{120.8}{pm} pada \ce{O2}, dan \qty{147.5}{pm} pada
        \ce{H2O2}. Apakah ramalan itu terbukti?
  \item Tuliskan tiga \omterm{def:b1:lewis-resonance-vsepr:resonance}{struktur resonansi} \ce{NO3-} dan tentukan orde
        setiap ikatan \ce{N-O}.
  \item Dalam \ce{HNO3} ketiga ikatan \ce{N-O} berukuran 140.6, 121.1,
        dan \qty{119.9}{pm}. Pasangkan nilai-nilai itu dengan ikatannya
        dan jelaskan dengan resonansi.
  \item Jelaskan mengapa sudut \ce{O-N-O} pada \ce{NO2} ($134^\circ$)
        lebih besar daripada $120^\circ$.
  \item Manakah di antara kedua struktur pada pertanyaan 4 yang
        menempatkan \omterm{def:b1:lewis-resonance-vsepr:formal-charge}{muatan formal} negatif pada atom yang lebih
        elektronegatif?
\end{enumerate}

\medskip
\noindent\textbf{Bagian III --- Bentuk dan kepolaran.}
\begin{enumerate}[resume]
  \item Tentukan tipe \ce{AX_nE_m} atom pusat \ce{O3}, \ce{N2O},
        \ce{NO3-}, dan nitrogen \ce{HNO3}, beserta bentuknya.
  \item Sudut ozon adalah $116.8^\circ$. Jelaskan mengapa sudut itu di
        bawah $120^\circ$.
  \item Ozon mempunyai \omterm{def:b1:lewis-resonance-vsepr:dipole}{momen dipol} \qty{0.53}{\debye}, \ce{N2O}
        \qty{0.16}{\debye}, \ce{N2} tidak ada. Jelaskan.
  \item Mengapa \ce{CO2} nonpolar sedangkan \ce{SO2} polar?
  \item Ramalkan sudut \ce{NH3} relatif terhadap $109.5^\circ$, lalu
        bandingkan dengan nilai terukur $106.7^\circ$.
  \item Pertanyaan yang sama untuk air ($104.5^\circ$).
\end{enumerate}
% ledger: geo:O3.rOO, re:O2, geo:H2O2.rOO, geo:HNO3.rNO, geo:HNO3.rNO2,
% geo:HNO3.rNO3, geo:NO2.aONO, geo:O3.aOOO, dip:O3, dip:N2O

\medskip
\noindent\textbf{Bagian IV --- Muatan parsial air.}
\omterm{def:b1:lewis-resonance-vsepr:dipole}{Momen dipol} air adalah \qty{1.857}{\debye}, sudutnya $104.48^\circ$, dan
panjang ikatan \ce{O-H}-nya \qty{95.8}{pm}.
% ledger: dip:H2O, geo:H2O.aHOH, geo:H2O.rOH
\begin{enumerate}[resume]
  \item Nyatakan \omterm{def:b1:lewis-resonance-vsepr:dipole}{momen dipol} air sebagai fungsi \omterm{def:b1:lewis-resonance-vsepr:dipole}{dipol ikatan} $\mu_{OH}$
        dan sudutnya.
  \item Hitunglah $\mu_{OH}$ dalam debye.
  \item Ubahlah nilai itu menjadi coulomb meter.
  \item Seandainya \omterm{def:b1:lewis-resonance-vsepr:dipole}{dipol ikatan} itu terbentuk dari dua muatan penuh
        $\pm e$ pada kedua inti, berapa nilainya?
  \item Turunkan sifat ion fraksional $\delta$ ikatan \ce{O-H}, yaitu
        rasio \omterm{def:b1:lewis-resonance-vsepr:dipole}{dipol ikatan} terukur terhadap dipol muatan penuh.
\end{enumerate}
\end{problem}
